% chapters/chapter02.tex
\chapter{Si sta facendo sera}

Caterina \`e impegnata nei preparativi per il suo volo transoceanico con cui spargerà dieci quintali di CO$_2$ nell'atmosfera. Giovanni attende con Ipparchia un po' di cibo recuperato dalla distribuzione ordinaria. La sua impronta di carbonio \`e bassa per necessità, non per scelta. Vive nell'ex convitto  "San Giorgio", senza i comfort di un'abitazione urbanizzata, e produrre pi\`u CO$_2$ di quella che espira \`e quasi impossibile.

\par\medskip
Ipparchia \`e seduta sotto una quercia vicino all’uscita sul retro, la \emph{dumpster zone}. Da un momento all’altro i commessi usciranno per le pulizie e quello sarà il segnale. Giovanni aspetta con calma, pronto a contrattare qualche prodotto in scadenza.\\

La porta si apre: Ipparchia capta il segnale e si fa avanti. Giovanni attende, non si muove subito: aspetta un cenno di invito. Ripensa alle parole di sua madre, un mantra: \textit{La prima regola  \`e la buona educazione}. Quando la ragazza  lo chiama, Giovanni ricambia con un sorriso. Lei rientra per uscire di nuovo con un cartone che non deve essere troppo pesante a giudicare dalla nonchalance con cui lo destreggia.

\begin{center}
  \includegraphics[width=.9\linewidth]{media/dumper_zone.png}
\end{center}



\par\medskip
Lo appoggia al pianerottolo delle scale di servizio. Ipparchia tira Giovanni e i due si avviano verso di lei, in equilibrio tra zelo e prudenza. Lei li aspetta  e mentre si avvicinano
 accende una sigaretta. \`E tanto giovane da credersi immortale.

\par\medskip
«Ciao Giò, ciao bella!»

\par\medskip
Ipparchia accetta la carezza e ricambia il saluto ricevuto.

\par\medskip
«Ti fa male fumare», le dice sorridendo.

Lei lo guarda per un istante. Tira ancora qualche boccata poi la spegne.

\par\medskip
«Torno dentro!»\\
«Ti ho rovinato la sigaretta.»\\
«Macché, devo riprendere il turno.» Apre il portone. «A presto, vi aspetto.»\\
«Grazie, prendo il pacco!»

\par\medskip
Adesso deve portare il pacco al convitto. La strada non \`e lunga, ma Giovanni deve ricordarla a memoria. Nella mente ha catalogato strade, stradelli, edifici, alberi e piante, organizzati in strati come nelle mappe  governative digitali.
\par\medskip
La vegetazione per\`o cresce rapidamente, l’edera inghiotte metri di strada ogni mese, i rami delle siepi invadono i marciapiedi, e la topografia ufficiale resta sempre un passo indietro rispetto al mondo vivo.

Ma tra i neuroni naturali di Giovanni si annidano connessioni che lo strato cognitivo del tecnocene non ha ancora registrato.

Lui ha esplorato siepi campestri, corridoi ecologici, fasce ripariali e margini di bosco. Quei luoghi sono l\`i, dentro di lui, con i loro merli, le capinere, i ricci, le arvicole, i fagiani, i ramarri e gli impollinatori. I nidi di quegli insetti non sono presenti nei database, ma il \textit{Prunus cerasifera} fruttifica anche grazie a loro, e Giovanni ne mangia chili ogni estate.

\par\medskip
Questa \`e la mappa che Giovanni ha nel cervello: odori, superfici, distanze misurate a passi e colpi di bastone.
Una  mappa che va aggiornata di continuo: dove la memoria vacilla, ci sono il fiuto e gli occhi di Ipparchia.


\begin{center}
  \includegraphics[width=.82\linewidth]{"media/mappa di odori.png"}
\end{center}


\par\medskip

\par\medskip
Ipparchia \`e al suo fianco da quando Giovanni \`e diventato un senza fissa dimora. Anche lei vagava sola e il loro incontro sarebbe stato degno di una nota nei rotocalchi se qualche giornalista lo avesse testimoniato. Questa \`e la vita di strada, capita di tutto, ma nessuno se ne cura. Ipparchia non lo tradirebbe mai.
\par\medskip

Un fagiano, disturbato dalla loro presenza, esce da una siepe e gli taglia la strada. Giovanni sente solo un fruscio veloce dell’erba, Ipparchia si lancia dietro al fagiano e i due scompaiono nel corridoio ecologico.
\par\medskip
Giovanni \`e solo, con uno scatolone in mano e la sua compagna da ritrovare prima che corra in un luogo fuori dalle sue mappe.

\par\medskip
Il silenzio si stende, pesante, tra le erbacce. Giovanni ascolta il suo stesso respiro, poi si rimette in cammino.




